\subsection{The Prym factor has rank zero and an exact two-primary Tate--Shafarevich group}

\begin{statement}{A.154}{Proved}
{An explicit Cassels--Tate calculation and the dyadic restriction on full Selmer sections}
Put
\[
 \begin{gathered}
 K=\Q(t),\quad t^2=2,\quad \epsilon=1+t,\\
 r=1119169,\quad M=325596479663,\\
 b=1146\epsilon^2M^2,\qquad E_b:\ y^2=x^3+bx.
 \end{gathered}
\tag{A.154.1}
\]
For the literal Prym factor and maps in A.150, one has
\[
 \begin{split}
 E_b(K)&=\{O,(0,0)\},\qquad \operatorname{rank}E_b(K)=0,\\
 \operatorname{Sel}^{2}(E_b/K)&\simeq(\Z/2\Z)^3,\\
 \Sha(E_b/K)[2^\infty]&=\Sha(E_b/K)[2]
                    \simeq(\Z/2\Z)^2,\qquad
 T_2\Sha(E_b/K)=0.
 \end{split}
\tag{A.154.2}
\]
The Cassels--Tate pairing on the displayed two-dimensional
Tate--Shafarevich group is nondegenerate.  These conclusions
use the exact six-class isogeny calculation below and an
explicit full two-cover pairing; no analytic rank assertion
or finiteness hypothesis on the full Tate--Shafarevich group
is used.

Let \(D^\circ\) be the particular harmonic curve of A.134,
equivalently the open part of
\[
 \overline D:\ aW^4=T(T-1)(T-\tfrac12),\qquad
 a=573\left(\frac{M}{2\cdot191^2}\right)^2,
\tag{A.154.3}
\]
obtained by deleting its four boundary points
\(P_0,P_1,P_{1/2},P_\infty\).
Every full Selmer section of \(D^\circ/\Q\) has cuspidal
localization at two.  The same assertion holds for a full
Selmer section of \(D^\circ_K/K\) at the unique place over
two.  More precisely, restriction, compactification, and
the actual positive-sign Prym map give a two-primary
elliptic section class equal to the global Kummer class of
one point \(Q\in\{O,(0,0)\}\).  Its dyadic proper source
point belongs to
\[
 \begin{cases}
 \{P_0,P_\infty\},&Q=O,\\
 \{P_1,P_{1/2}\},&Q=(0,0).
 \end{cases}
\tag{A.154.4}
\]
This supplies a restriction at the stated place; the
proof below also identifies the ordinary points at odd
good places which have zero two-primary Kummer class.
\end{statement}

\begin{proof}
\emph{The independent full-Selmer upper bound.}
We first prove the order bound using the actual two-isogenies
and the already determined groups of A.150.  The following
calculation is included with its divisor maps and local
certificates.

\paragraph{A finite second-descent bound from the six isogeny classes.}
Write \(A=E_b\), \(B=E_{-4b}\), and let
\(\phi:A\to B\), \(\psi:B\to A\) be the actual dual maps
of A.150.16.  In particular \(\psi\phi=[2]_A\) and
\(\phi\psi=[2]_B\).  We prove
\[
                      |\operatorname{Sel}^{2}(A/K)|\leq8
\tag{A.154.S1}
\]
using only the two already proved groups in A.150.6
and the explicit finite calculation below.

For \(d\in\mathcal S_b\), put \(B_d=b/d\) and
\[
 C_d:\ N^2=dU^4+B_dV^4,\qquad
 \rho_d:C_d\longrightarrow A,\quad
 (U:V:N)\longmapsto
 \left(\frac{dU^2}{V^2},\frac{dUN}{V^3}\right).
\tag{A.154.S2}
\]
As proved by the actual inverse formula in A.150.18--20,
over \(K(\sqrt d)\) this is the map \(\psi:B\to A\).
Changing the sign of \(\sqrt d\) translates the source
by \(T'=(0,0)\in B[\psi]\).  Thus
\[
 \mathcal S_b=\operatorname{Sel}^{\psi}(B/K)
     \longrightarrow \Sha(B/K)[\psi],\qquad
 d\longmapsto [C_d],
\tag{A.154.S3}
\]
is induced by the literal inclusion
\(B[\psi]=\{O,T'\}\hookrightarrow B\):
the cocycle of \(C_d\) is \(g\mapsto\kappa_d(g)T'\).
This identifies its Jacobian with \(B\), including the
specified kernel point.  There is no unspecified
identification of genus-one models here.

\emph{The tangent function and its divisor class.}
Choose \(n_d,z_d,v_d\in K\) with
\[
                  n_d^2=dz_d^2+B_dv_d^2 .
\tag{A.154.S4}
\]
The six choices used below have \(v_d=1\), and all
three coordinates are nonzero.  The map
\[
 C_d\longrightarrow Q_d,\qquad
  (U:V:N)\longmapsto(N:U^2:V^2),\qquad
 Q_d:\ n^2=dz^2+B_dv^2
\tag{A.154.S5}
\]
is the degree-two quotient by
\(j(U:V:N)=(-U:V:N)\).
The line tangent to \(Q_d\) at \((n_d:z_d:v_d)\)
pulls back to
\[
 \ell_d=n_dN-dz_dU^2-B_dv_dV^2.
\]
The exact polynomial identity is
\[
 \ell_d\bigl(n_dN+dz_dU^2+B_dv_dV^2\bigr)
          =b(v_dU^2-z_dV^2)^2.
\tag{A.154.S6}
\]
Its geometric divisor is twice the degree-two fibre
of (A.154.S5).  Consequently
\[
 f_d=\ell_d/V^2,\qquad
       \operatorname{div}(f_d)=2D_d,\qquad
       D_d=F_j-F_\iota\in\operatorname{Div}^0(C_d),
\tag{A.154.S7}
\]
where \(F_j\) is that conic fibre and \(F_\iota\)
is the fibre at infinity of the hyperelliptic map
\((U:V:N)\mapsto(U:V)\).
Its involution is \(\iota(U:V:N)=(U:V:-N)\).
For a degree-two map from a genus-one curve the
involution is \(P\mapsto F-P\) in the corresponding
degree-two divisor class.  Thus \(j\iota\) is
translation by \([F_j-F_\iota]\).
The literal formula
\[
                    j\iota(U:V:N)=(-U:V:-N)
\]
fixes both functions in (A.154.S2), and is its
nonidentity deck transformation.  Under its already
specified Jacobian \(B\), it is translation by \(T'\).
Therefore \([D_d]=T'\in\operatorname{Pic}^0(C_d)(K)\).
This is also the divisor calculation of
\cite[Lemma 6.1]{AgentFisherBinaryQuartics2022}.

For \(e\in\mathcal S_b\), the quaternion
\((e,f_d)\) is unramified on the proper curve,
because (A.154.S7) has even divisor.
Over \(K(\sqrt e)\), the norm of the rational
divisor \(D_d\) is \(2D_d=\operatorname{div}(f_d)\).
The connecting map from divisors to the Brauer group
therefore sends the cocycle \(\kappa_e T'\) to
\((e,f_d)\).  This proves, with its actual cocycle
and function, the formula
\[
 (-1)^{\,2\langle[C_d],[C_e]\rangle_{\rm CT}}
       =\prod_w(e,f_d(P_w))_w,\qquad P_w\in C_d(K_w).
\tag{A.154.S8}
\]
Here the exponent means the identification
\(\tfrac12\Z/\Z\simeq\{\pm1\}\).
Equivalently the additive invariant is the sum
of the quaternion invariants.
This is the homogeneous-space definition and
cyclic-algebra map of
\cite[Section 4, equation (10), and Lemma 5.1]
{AgentFisherBinaryQuartics2022}.
At a chart with \(V=0\), use \(\ell_d/U^2\);
the two functions differ by a square wherever
both charts are defined, so this is the same
Brauer evaluation.  This accounts for every
chart used in the finite certificate.

\emph{Six literal conic points and the finite calculation.}
Let \(d_1,\ldots,d_6\) be the ordered basis of
\(\mathcal S_b\) in A.150.6, let \(t=\sqrt2\),
and retain \(M\) from A.150.1.  The following table
gives \(n_{d_i}=M\nu_i\), \(z_{d_i}=M\zeta_i\),
\(v_{d_i}=1\).  Direct substitution verifies
\(\nu_i^2=d_i\zeta_i^2+1146(1+t)^2/d_i\).
\[
 \begin{array}{c|c|c}
 i&\nu_i&\zeta_i\\ \hline
 1&-50-26t&-18-17t\\
 2&-26(1+t)/7&(2+t)/7\\
 3&(90+37t)/7&(-262-78t)/161\\
 4&(182-163t)/21&(10596-8182t)/67137\\
 5&(-47-29t)/17&(4098363+2550620t)/133181111\\
 6&(6196+4761t)/3361&
           (584427298+410557178t)/26330689063
 \end{array}
\tag{A.154.S9}
\]
These points were found by a conic search; their
displayed equations, rather than an output about
class groups or Selmer groups, are the entire
input from that search.

For each row, multiply the coefficients of
\(\ell_d\) by a common denominator and remove
their greatest common divisor in
\(\mathcal O_K=\Z[t]\).
The Euclidean algorithm of A.150 proves that
this is legitimate.  The certificate retains
explicit integral Bezout coefficients showing
that the three resulting coefficients generate
the unit ideal.  Multiplying \(f_d\) by this
single global constant does not change the
product in (A.154.S8), by global quaternion
reciprocity.

The only rational finite primes used are
\[
             2,3,5,7,13,23,139,191,1119169.
\tag{A.154.S10}
\]
Every \(d_i\) and \(b/d_i\) is integral and
supported on this list; this follows either
from A.150.22--23 or literal division in
\(\Z[t]\).  The second symbol entries in
(A.154.S8) use these \(d_i\), so they are units
outside this list.
At every other prime the quartic has smooth
genus-one reduction.  The primitive weighted
quadratic tangent cannot vanish identically
on that curve: its reduction has nonzero
coefficients and weighted degree two, whereas
the defining equation has weighted degree
four.  Its zero divisor has degree four.
The residue field has size \(q\geq11\), and
\[
                   q+1-2\sqrt q>4.
\]
There is consequently a residue point at
which the tangent is a unit.  One of \(U,V\)
is a unit; divide by its square and lift the
smooth point.  The resulting Hilbert symbol
of two units at an odd place is one.
This proves the vanishing at every omitted
finite place.  All \(d_i\) are totally positive,
so the two real symbols are also one.

Here are the exact local operations in the
retained certificate.  At a split odd prime
the residue map is \(t\mapsto t_w\), with its
unique Hensel lift from \(t_w^2=2\); at an inert
prime it is the actual field
\(\F_p[t]/(t^2-2)\).
For \(a=\varpi^\alpha a_0\), \(c=\varpi^\gamma c_0\)
over an odd residue field \(\F_q\), it uses
\[
 (a,c)_w=(-1)^{\alpha\gamma(q-1)/2}
              \overline a_0^{\,\gamma(q-1)/2}
              \overline c_0^{\,\alpha(q-1)/2}.
\tag{A.154.S11}
\]
This is the elementary quadratic norm test:
the unramified quadratic norm group consists
of the elements of even valuation; the
ramified quadratic norm on units is the
residue-square test.  Bilinearity and
\((a,-a)_w=1\) give the displayed formula.
For an inert place, the quadratic character
of a residue \(a+ct\) equals the Legendre
symbol of \(a^2-2c^2\) in \(\F_p\), since
this is the residue norm.

At the dyadic place use the squareclass basis
\((-1,5,\epsilon,\pi)\), where
\(\epsilon=1+\pi\) and \(\pi=\sqrt2\).
In additive \(\F_2\)-notation its Hilbert matrix is
\[
 \begin{pmatrix}
 0&0&1&1\\
 0&0&0&1\\
 1&0&1&1\\
 1&1&1&1
 \end{pmatrix}.
\tag{A.154.S12}
\]
For clarity, this matrix has a direct norm proof.
Two rational entries pair trivially after the
degree-two extension from \(\Q_2\).
Corestriction and
\(N_{K_2/\Q_2}(\epsilon)=-1\),
\(N_{K_2/\Q_2}(\pi)=-2\) give
\((\epsilon,-1)=(\pi,-1)=-1\) and
\((\epsilon,5)=1\), \((\pi,5)=-1\).
Here \(5\) defines the unramified quadratic
extension of \(\Q_2\), whose norms have even
valuation.  The remaining rational test
\((-1,-1)_{\Q_2}=-1\) follows since a sum of
two squares cannot be minus one in \(\Q_2\);
\((2,-1)_{\Q_2}=1\) follows from
\(2=N_{\Q_2(i)/\Q_2}(1+i)\).
The Steinberg relation
\((\epsilon,1-\epsilon)=(\epsilon,-\pi)=1\)
gives \((\epsilon,\pi)=-1\), and
\((a,a)=(a,-1)\) gives the two diagonal
entries.  These are every entry of (A.154.S12).
The basis is exhaustive: its sixteen products
give all sixteen squareclasses enumerated
modulo \(\pi^5\) in A.150.25.

All chosen local coordinates are retained
as exact rational pairs in \(K\).
For an ordinary odd square-root lift,
an approximation whose squared error has
valuation at least \(L+v(f)\) differs from
the actual root by valuation at least
\(L+v(f)/2\).
For an odd fourth-root lift the derivative
is a unit.  At two the certificate divides
out an even valuation, starts from a unit
square root modulo \(\pi^5\), and performs
Newton steps; its error bound on the root
is \(L+v_\pi(f)/2-2\).
Every tangent evaluation is checked stable
by a margin at least one at odd places and
five at two.  Thus the recorded signs are
symbols at actual local points, rather than
at unlifted residue points.

The resulting additive matrix on the
ordered basis \(d_1,\ldots,d_6\) is
\[
 \mathcal M=
 \begin{pmatrix}
 0&0&0&1&1&1\\
 0&0&0&1&1&1\\
 0&0&0&1&0&0\\
 1&1&1&0&0&1\\
 1&1&0&0&0&0\\
 1&1&0&1&0&0
 \end{pmatrix},
 \qquad \operatorname{rank}_{\F_2}\mathcal M=4 .
\tag{A.154.S13}
\]
Its radical is explicitly
\(\langle d_1d_2,d_3d_5d_6\rangle
 =\langle573,23\cdot1119169\rangle\):
the second product is literally
\(23\cdot7^2\cdot1119169\).
The literal row operations, six conic equations,
primitive Bezout identities, and every local
witness and sign are executed in
\path{code/harmonic_prym_isogeny_pairing_certificate.py}
and its retained output and data.

\emph{The exact Selmer maps giving the bound.}
The finite group schemes fit into
\[
 0\longrightarrow A[\phi]\longrightarrow A[2]
       \xrightarrow{\phi}B[\psi]\longrightarrow0.
\tag{A.154.S14}
\]
Its cohomology maps, with the local Kummer
conditions, give
\[
 \operatorname{Sel}^{\phi}(A/K)
       \longrightarrow\operatorname{Sel}^{2}(A/K)
       \xrightarrow{\phi_*}\operatorname{Sel}^{\psi}(B/K).
\tag{A.154.S15}
\]
The kernel of the first map is the boundary
of \(B[\psi](K)/\phi A[2](K)\).
Both rational kernels have order two,
\(\phi A[2](K)=0\), and this boundary is
the nonzero class \([-4b]\).
Since A.150 gives
\(|\operatorname{Sel}^{\phi}(A/K)|=4\),
its image has order two.
This is the whole kernel of the second
map.  To check the local condition in
that exactness assertion, take a global
lift in \(H^1(K,A[\phi])\) of a Selmer
class whose second image is zero.
At \(w\), write the original class
as \(\delta_2(P_w)\).
Its zero \(\psi\)-Kummer image says
\(P_w=\psi Q_w\) for \(Q_w\in B(K_w)\).
Subtracting \(\delta_\phi(Q_w)\) leaves
the boundary of \(B[\psi](K_w)\), which
is itself a local \(\phi\)-Kummer class.
Thus the global lift satisfies every
required local condition.

For \(\xi\in\operatorname{Sel}^{2}(A/K)\),
let \(\bar\xi\in\Sha(A/K)[2]\) be its
image.  The torsor image of \(\phi_*\xi\)
under (A.154.S3) is literally
\(\phi_*\bar\xi\in\Sha(B/K)\).
For every \([C_e]\) from (A.154.S3),
\(\psi_*[C_e]=0\).  The adjointness of
the actual dual isogenies therefore gives
\[
 \langle\phi_*\bar\xi,[C_e]\rangle_B
       =\langle\bar\xi,\psi_*[C_e]\rangle_A=0.
\tag{A.154.S16}
\]
The adjointness statement and its proof
by pulling back the divisor functions are
\cite[Theorem 1.2 and Section 2, pp.\ 187--188]
{AgentCasselsIsogenies1965}.
Thus \(\phi_*\operatorname{Sel}^{2}(A/K)\)
lies in the radical of the proved matrix
(A.154.S13), which has dimension two.
Combining this with the kernel of order
two in (A.154.S15) proves (A.154.S1).
No order of a computed full Selmer group,
and no finiteness hypothesis on a
Tate--Shafarevich group, enters this bound.


\paragraph{A nonzero pairing on full two-covers.}
We now exhibit two everywhere locally soluble two-covers of
\(E_b\) whose Cassels--Tate pairing is \(1/2\).
All subsequent data are literal rational numbers in
\(K=\Q(t)\).  In particular, the order of a full Selmer
group printed by the discovery software is not an input
to this proof.

For a binary quartic
\(g(X,Z)=AX^4+BX^3Z+CX^2Z^2+DXZ^3+EZ^4\), put
\[
 \begin{split}
 I(g)&=12AE-3BD+C^2,\\
 J(g)&=72ACE-27AD^2-27B^2E+9BCD-2C^3.
 \end{split}
\tag{A.154.5}
\]
The three quartics used below all have
\(I=-48b\), \(J=0\).  Their discriminant
\(16(4I^3-J^2)/27\) is nonzero.
The associated elliptic curve in Fisher's convention is
\(E_{I,0}:Y^2=X^3-27IX\).  The actual isomorphism is
\[
 E_b\longrightarrow E_{I,0},\qquad
               (x,y)\longmapsto(36x,216y).
\tag{A.154.6}
\]

Here is also the covering morphism, to fix its model and
normalization explicitly.  Write \(f(u)=g(u,1)\), and set
\[
 H_f=3(f')^2-4ff'',\qquad
 R_f=4(f'H_f-fH_f').
\tag{A.154.7}
\]
Direct polynomial expansion gives
\[
             3R_f^2-16H_f^3-36864bH_ff^2=0.
\tag{A.154.8}
\]
Thus on the smooth projective completion of \(N^2=f(u)\)
the rational functions
\[
 q_g(u,N)=
       \left(\frac{H_f(u)}{48N^2},
             \frac{R_f(u)}{768N^3}\right)
\tag{A.154.9}
\]
define a morphism to \(E_b\).  They extend to the proper
curve because the target is proper and the source is
smooth.  At each of the four distinct geometric branch
points \(N=0\), \(H_f=3(f')^2\ne0\); the pole orders of
these functions are two and three.  Each branch point
therefore maps to \(O\), with ramification index one.
There are no other poles, including the two points at
infinity when \(g(1,0)\ne0\).  Hence the morphism has
degree four and is etale by Riemann--Hurwitz.
Choose a geometric branch point as origin.  The other
three differ from it by the three nonzero two-torsion
points, since the hyperelliptic involution is negation
with this origin.  Consequently the morphism is an actual
two-cover: it is multiplication by two followed by the
specified isomorphism of its Jacobian with \(E_b\).
This factorization descends to \(K\): the induced Jacobian
homomorphism is defined over \(K\), and its factor through
\([2]\) is unique because the homomorphism group of two
elliptic curves has no nonzero torsion.
The differential identity
\[
                  q_g^*\frac{dx}{2y}=-\frac{2\,du}{N}
\tag{A.154.10}
\]
and (A.154.6) give the normalization used here.
The corresponding quartic description of the finite
Kummer class is exactly
\cite[Lemma 2.1 and Section 2]{AgentFisherBinaryQuartics2022}.

Use
\[
 g_1(X,Z)=r(X+Z)^4-\frac{4b}{r}X^4.
\tag{A.154.11}
\]
The change \((X,Z)=(u,1-u)\), of determinant one, gives
the small quartic
\[
                  C_r:\ N^2=r-\frac{4b}{r}u^4.
\tag{A.154.12}
\]
This curve is everywhere locally soluble by the exact
class \(r\in\mathcal S_{-4b}\) of A.150.  Its direct map
to \(E_{-4b}\) is
\((u,N)\mapsto(r/u^2,rN/u^3)\); composing with \(\psi_b\)
gives
\[
 x_0=\frac{N^2}{4u^2},\qquad
 y_0=\frac{N(r^2+4bu^4)}{8ru^3}.
\tag{A.154.13}
\]
The covariant map (A.154.9), expressed on (A.154.12), is
the translate of this map by \(T_0=(0,0)\):
\[
 (x_0,y_0)\longmapsto
           \left(\frac b{x_0},-\frac{by_0}{x_0^2}\right).
\tag{A.154.14}
\]
Indeed \(H_f=192bu^2\) and
\(R_f=-1536bu(r+4bu^4/r)\).
This records the possible torsion translation explicitly;
all uses below concern the displayed quartic normalization.

\paragraph{The literal Fisher input coefficients.}
Here \(t=\sqrt2\), \(r=1119169\), and
\(g_i(X,Z)=\sum_{j=0}^4 c_{i,j}X^jZ^{4-j}\).
The first polynomial is
\(g_1=r(X+Z)^4-(4b/r)X^4\).
The remaining coefficients, and the chosen square root
\(m=m_0+m_1\varphi+m_2\varphi^2\), are as follows.
Every displayed integer is an exact input, and the same
rational pairs are retained in
\path{code/harmonic_prym_fisher_exact_data.json}.

\[
c_{2,0}=\frac{\begin{gathered}\scriptstyle -1105588195179791625651078095706697333069\\[-2pt]\scriptstyle {}+877556745246639942323602265475930377568t\end{gathered}}{935555900741191077811391382507783188961226470515575}.
\tag{A.154.D20}
\]

\[
c_{2,1}=\frac{\begin{gathered}\scriptstyle -5540928743512790350006981292953388\\[-2pt]\scriptstyle {}+9348210277430937552637225227327984t\end{gathered}}{1087607601378646157938938703818175}.
\tag{A.154.D21}
\]

\[
c_{2,2}=\frac{\begin{gathered}\scriptstyle 71783834352841392547954860666540\\[-2pt]\scriptstyle {}+120955658328769591100673151637112t\end{gathered}}{2284719449267034025}.
\tag{A.154.D22}
\]

\[
c_{2,3}=\frac{\begin{gathered}\scriptstyle 1757082965015780717118550076480888505038563656\\[-2pt]\scriptstyle {}+1334562270654411275381324637376168444287220368t\end{gathered}}{2284719449267034025}.
\tag{A.154.D23}
\]

\[
c_{2,4}=\frac{\begin{gathered}\scriptstyle -7561243441172084609148446045366886357854699902448831772732\\[-2pt]\scriptstyle {}-5343266287278555475164967326114882400164713497381811727024t\end{gathered}}{2284719449267034025}.
\tag{A.154.D24}
\]

\[
c_{3,0}=\frac{\begin{gathered}\scriptstyle 41455951009817066885052607588499608177788939499\\[-2pt]\scriptstyle {}-27768599768386499529560925513889558238916961380t\end{gathered}}{2125947929722471429942234964145147052205608530886862085300}.
\tag{A.154.D30}
\]

\[
c_{3,1}=\frac{\begin{gathered}\scriptstyle 101278412967185326494385306629415677088532\\[-2pt]\scriptstyle {}-52095764030506724995200175292601444292679t\end{gathered}}{2700261657400188170244993859108806176385}.
\tag{A.154.D31}
\]

\[
c_{3,2}=\frac{\begin{gathered}\scriptstyle 34665317701806075883725801409853605497\\[-2pt]\scriptstyle {}-1088923055185222909612383981944999312t\end{gathered}}{873436184739855072128521}.
\tag{A.154.D32}
\]

\[
c_{3,3}=\frac{\begin{gathered}\scriptstyle 49143251797310725165795561681745556654464062765400\\[-2pt]\scriptstyle {}+27360957632102280199593493880778823962642438265610t\end{gathered}}{873436184739855072128521}.
\tag{A.154.D33}
\]

\[
c_{3,4}=\frac{\begin{gathered}\scriptstyle -1145516250448259949120335174735880752753369070770615441241654825\\[-2pt]\scriptstyle {}-809469008979191424190242184343438826054550107506776656562735500t\end{gathered}}{873436184739855072128521}.
\tag{A.154.D34}
\]

\[
m_0=\frac{\begin{gathered}\scriptstyle -2230575953574473918670466577224295377555585357344976900617784771146240\\[-2pt]\scriptstyle {}-1577671320202015910344669780113251001931088624068196254982397769760256t\end{gathered}}{19}.
\tag{A.154.Dm0}
\]

\[
m_1=\frac{\begin{gathered}\scriptstyle -1112580307357952746036636417589868462854274589857433525437910236391657104660672\\[-2pt]\scriptstyle {}-786525733221423318336821760821550148588617948894045868365395216213846024996544t\end{gathered}}{282528344699694941509}.
\tag{A.154.Dm1}
\]

\[
m_2=\frac{\begin{gathered}\scriptstyle 1710872698619339936174938673988957021297294051351083748147167896\\[-2pt]\scriptstyle {}+1217260528857972263698557756110138500060072782019046995324168100t\end{gathered}}{847585034099084824527}.
\tag{A.154.Dm2}
\]


\emph{Everywhere local solubility of the other two covers.}
Let \(g_2,g_3\) have the preceding literal coefficients.
Substitution in (A.154.5) gives \(I=-48b,J=0\) for both.
The rational primes containing every coefficient denominator
and every prime dividing \(2b\) are respectively
\[
 \begin{split}
 S_2={}&\{2,3,5,7,13,23,31,139,191,2857,6473,r,9751799\},\\
 S_3={}&\{2,3,5,7,13,19,23,79,139,191,211,r,622636961\}.
 \end{split}
\tag{A.154.15}
\]
At every place over these sets, the default affine coordinate
is \(u=0\), except for the entries in the following table.
A listed residue root \(t_0\) means the place with
\(t\equiv t_0\pmod p\); a dash means the unique place.
An infinity entry means that the leading coefficient is
a nonzero square in the local field.
\[
\begin{array}{c|r|r|c}
\text{curve}&p&t_0&u\\\hline
g_2&2&-&\infty\\
g_2&5&-&1+4t\\
g_2&13&-&t/13\\
g_2&31&8&\infty\\
g_2&31&23&9\\
g_2&139&-&(66+130t)/139\\
g_2&2857&2657&3\\
g_2&6473&3861&3\\
g_2&9751799&2503412&2868851\\\hline
g_3&2&-&-8-15t/2\\
g_3&3&-&(1+t)/3\\
g_3&19&-&2+18t\\
g_3&79&9&70\\
g_3&191&57,134&\infty\\
g_3&622636961&449895035&408991746
\end{array}
\tag{A.154.16}
\]
Every evaluation in this table and at the default points
is nonzero.  At an odd place its valuation is even and
its normalized residue is a square.  Hensel's lemma for
the unit square root therefore gives an actual point.
At two, reduction of the unit part modulo \(\pi^5\)
gives one of the two squares \(1,3+2\pi\) of A.150,
with even valuation.  If a unit has that square residue,
a lift of the residue root has error valuation at least
five and derivative valuation two; the strict inequality
\(5>2\cdot2\) gives an actual square root.  This proves
the dyadic assertions as well.
At the two real embeddings, \(g_2(0,1)>0\) at \(t>0\)
and \(g_2(-1/M,1)>0\) at \(t<0\); \(g_3(0,1)>0\)
at both.  These inequalities are checked by comparing
the squares of rational coefficients of \(a+ct\).

Outside (A.154.15), the integral binary quartic has unit
discriminant, so its weighted projective equation is a
smooth genus-one curve over the residue field.  The
genus-one bound \(\#C(\F_q)\geq q+1-2\sqrt q>0\)
provides a residue point, which lifts by smoothness.
Thus both curves are everywhere locally soluble.
All exceptional evaluations, including their normalized
residue characters and the denominator-support checks,
are executed in
\path{code/harmonic_prym_fisher_els_certificate.py}.

\emph{The exact cubic algebra and pairing function.}
Use the algebra
\[
 L_I=K[\varphi]/(\varphi^3-3I\varphi),\qquad
 z(g)=\frac{4A\varphi+3B^2-8AC}{3}.
\tag{A.154.17}
\]
In the original elliptic splitting algebra
\(K[\theta]/(\theta^3+b\theta)\), the actual algebra
isomorphism is \(\varphi\mapsto-12\theta\).
The coefficients \(m_i\) displayed above satisfy the
literal equality
\[
       (m_0+m_1\varphi+m_2\varphi^2)^2
                         =z(g_1)z(g_2)z(g_3)\in L_I^\times.
\tag{A.154.18}
\]
In particular the three Selmer classes sum to zero under
the exact cubic-invariant group law in Fisher's Section 2.

For the Hessian binary quartic
\[
\begin{split}
h_g={}&(3B^2-8AC)X^4+4(BC-6AD)X^3Z\\
 &+2(2C^2-24AE-3BD)X^2Z^2\\
 &+4(CD-6BE)XZ^3+(3D^2-8CE)Z^4,
\end{split}
\tag{A.154.19}
\]
put
\[
 \mathcal G_g=(4\varphi g+h_g)/3,\qquad
 \mathcal H_g=\tfrac1{12}\partial_X^2\mathcal G_g
                         +\tfrac29(I-\varphi^2)Z^2.
\tag{A.154.20}
\]
Take the coefficient of \(\varphi^2\) in
\(z(g_2)z(g_3)\mathcal H_{g_1}/m\), then substitute
\((X,Z)=(u,1-u)\).  Dividing this quadratic by one
nonzero global constant gives the primitive integral form
\[
 \begin{split}
 \Gamma(u)={}&66876882465209-47288996359904t\\
 &+(-328789960652106096+232490058105731322t)u\\
 &+(398753556190578319224-281962130652772446708t)u^2.
 \end{split}
\tag{A.154.21}
\]
The certificate retains both the original form and the
removed constant.  It also gives integral Bezout coefficients
whose dot product with the three coefficients in
(A.154.21) is one, proving primitivity at every finite
place, including split places.  A global constant changes
the local symbol product by one by Hilbert reciprocity.
The other Hilbert entry is exactly
\[
 h=g_2(1,0)=
 \frac{\begin{gathered}\scriptstyle
 -7561243441172084609148446045366886357854699902448831772732\\[-2pt]
 \scriptstyle{}-5343266287278555475164967326114882400164713497381811727024t
 \end{gathered}}{2284719449267034025}.
\tag{A.154.22}
\]
Fisher's exact formula
\cite[Theorem 3.1 and Remark 3.2]{AgentFisherBinaryQuartics2022}
therefore gives
\[
       (-1)^{2\langle[g_1],[g_2]\rangle_{\rm CT}}
                 =\prod_w(h,\Gamma(u_w))_w,
       \qquad N_w^2=r-4bu_w^4/r.
\tag{A.154.23}
\]

\emph{The complete local product.}
The finite rational support used in (A.154.23) is
\[
 S=\{2,3,5,7,13,23,31,41,139,191,2857,6473,
                           9473,r,9751799,38709401\}.
\tag{A.154.24}
\]
This includes the separate denominators of both rational
coefficients of \(h\), all prime factors of its nonzero
norm, and the discriminant support of (A.154.12).
Taking only the norm denominator would not establish
integrality at both primes over a split rational prime;
the separate denominator check is part of the certificate.
At every place over \(S\) use \(u_w=0\), except
\[
\begin{array}{r|r|r}
p&t_0&u_w\\\hline
9473&1373,8100&2\\
r&230386&150136581454\\
r&888783&228527708220\\
38709401&15703973&1\\
38709401&23005428&2.
\end{array}
\tag{A.154.25}
\]
Every resulting right side of (A.154.12) is a nonzero
local square.  At the two places over \(r\) it has
valuation two; at the other odd places it is a square
unit.  At two, \(r\equiv1\pmod8\) gives the point
with \(u=0\).  At both real places the same coordinate
gives the positive value \(r\).

The only nontrivial symbols in the resulting complete
table are
\[
\begin{array}{r|r|c|c|c}
p&t_0&v_w(h)&
 (\overline{\Gamma(u_w)}\mid\F_w)&(h,\Gamma(u_w))_w\\\hline
41&17&1&-1&-1\\
139&-&3&-1&-1\\
38709401&15703973&1&-1&-1.
\end{array}
\tag{A.154.26}
\]
Here \(\Gamma(u_w)\) is a unit in all three rows.
The residue field at 139 is the actual quadratic field
\(\F_{139}[t]/(t^2-2)\), and its quadratic character is
computed by the norm to \(\F_{139}\).
At every other odd place over \(S\), the valuation and
residue-square formula (A.154.S11) gives one.
At two both \(h\) and \(\Gamma(0)\) are squares; the
literal unit-square test modulo \(\pi^5\) verifies this.
At both real places \(h<0\) and \(\Gamma(0)>0\), so
those symbols are also one.

For every omitted finite place, \(h\) is a unit and
(A.154.12) has smooth genus-one reduction.  The
homogenization of the primitive quadratic \(\Gamma\)
has at most four zeros on that reduction, and there
are at most two points at infinity.  The smallest
possible residue size outside \(S\) is seventeen:
the omitted prime eleven is inert and has residue
size \(121\).  Thus \(q+1-2\sqrt q>6\).
There is an affine smooth residue point at which
\(\Gamma\ne0\); its lift has both symbol entries
units.  Its symbol is one.
This is the exact finite reduction in
\cite[Remark 3.3]{AgentFisherBinaryQuartics2022}.
The independence of the product from the selected local
points follows from the homogeneous-space definition,
\cite[Section 4, equation (10)]{AgentFisherBinaryQuartics2022}:
the second locally trivial torsor makes the associated
Brauer class locally constant on the first curve.
Thus (A.154.26) proves
\[
                \langle[g_1],[g_2]\rangle_{\rm CT}=\tfrac12.
\tag{A.154.27}
\]

All identities (A.154.5)--(A.154.27), every finite field
embedding and local square test, the exact coefficient
support, and the 27 finite local evaluations are retained
in the three independently executable certificates
\path{code/harmonic_prym_fisher_maps_certificate.py},
\path{code/harmonic_prym_fisher_local_certificate.py}, and
\path{code/harmonic_prym_fisher_els_certificate.py}.
They use exact rational arithmetic and finite residue
calculations.  The displayed quartics were found with
Magma; their equations and the literal certificates are
the arithmetic inputs to this proof.

\paragraph{Rank, the full two-primary group, and the section consequence.}
The Cassels--Tate pairing is bilinear and alternating,
and its pullback to a Selmer group vanishes on the
Mordell--Weil Kummer image
\cite[Introduction, Remark 3.2(v), and Section 4]
{AgentFisherBinaryQuartics2022}.
By A.152, the nonzero rational two-torsion point \(T_0\)
has no rational half: a half would be a point of order
four.  Its image in \(E_b(K)/2E_b(K)\) is therefore
nonzero and lies in the radical of the Selmer pairing.
The nonzero pairing (A.154.27) supplies a two-dimensional
nondegenerate subspace.  Together with (A.154.S1), this
forces the Selmer group to have dimension exactly three
and its radical to have dimension exactly one.
Finite generation of \(E_b(K)\), in the precise finite
basis form used in A.150, now gives
\(\dim_{\F_2}E_b(K)/2E_b(K)=1+\operatorname{rank}E_b(K)\).
Its Kummer image lies in that one-dimensional radical,
so the rank is zero.  A.152 then gives the first line
of (A.154.2).

The finite Kummer exact sequence yields
\[
 0\longrightarrow E_b(K)/2E_b(K)
 \longrightarrow\operatorname{Sel}^{2}(E_b/K)
 \longrightarrow\Sha(E_b/K)[2]\longrightarrow0.
\tag{A.154.28}
\]
Hence the last group has dimension two and its induced
pairing is nondegenerate.  If \(z\in\Sha(E_b/K)[4]\),
then, for every \(w\in\Sha(E_b/K)[2]\), bilinearity gives
\[
                   \langle2z,w\rangle
                              =\langle z,2w\rangle=0.
\tag{A.154.29}
\]
Nondegeneracy on the two-torsion forces \(2z=0\).
An element of order \(2^n\), \(n>2\), would have a
multiple of order four, so none exists.  This proves
the exact equality \(\Sha[2^\infty]=\Sha[2]\), without
assuming finiteness of any larger group.  Since the
transition maps defining \(T_2\Sha\) are multiplication
by two, they are now zero; thus \(T_2\Sha=0\).

Let \(s\) be a full Selmer section as in the statement.
Restriction to \(G_K\), the open immersion into the proper curve, and
the actual map of A.150 define
\[
 s\longmapsto s_K\longmapsto\overline s_K
 \longmapsto(q_+)_*\overline s_K
 \longmapsto\eta\in H^1_{\rm cts}(G_K,T_2E_b).
\tag{A.154.30}
\]
Here the target extension is the literal
\(T_2E_b\rtimes G_K\) split by \(O\), as proved using
the multiplication covers in A.152.17.
At a local ordinary point the composite is its image
under \(q_+\); at a cusp compactification kills the
boundary inertia and gives the corresponding proper
boundary point.  These maps commute with every local
restriction, including the real places.

Project \(\eta\) to the finite quotient \(E_b[2^n]\),
and denote its class by \(\eta_n\).
Every localization is the finite Kummer class of the
actual proper elliptic image just described.  Therefore
\(\eta_n\in\operatorname{Sel}^{2^n}(E_b/K)\), and these
classes form a compatible inverse system.  The finite
multiplication sequences give
\[
 0\longrightarrow E_b(K)/2^nE_b(K)
 \longrightarrow\operatorname{Sel}^{2^n}(E_b/K)
 \longrightarrow\Sha(E_b/K)[2^n]\longrightarrow0.
\tag{A.154.31}
\]
The transition on the right is multiplication by two;
on the left it is the ordinary quotient map.
The left groups are the same finite group \(\{O,T_0\}\)
at every level, with identity transition.  Taking
inverse limits is exact here: compatible lifts are
obtained from the nonempty finite sets of choices at
each finite collection of levels, by compactness of
their product.  Thus
\[
 \varprojlim_n\operatorname{Sel}^{2^n}(E_b/K)
                      =\kappa_2\bigl(\{O,T_0\}\bigr).
\tag{A.154.32}
\]
This also follows directly: the compatible Sha images
are zero because \(T_2\Sha=0\), and uniqueness in the
left injections of (A.154.31) makes all the resulting
torsion points equal.

For completeness, no ambiguity from continuous
cohomology is hidden in this comparison.  If two
continuous \(T_2E_b\)-cocycles agree in cohomology at
every finite level, the finite sets of zero-cochains
which express their difference as a coboundary have
compatible nonempty solutions on each finite string
of levels.  Compactness gives a compatible sequence
of zero-cochains, hence a single element of \(T_2E_b\)
whose coboundary is the difference.  The map from
continuous cohomology to the inverse limit of these
finite cohomology groups is therefore injective in
the situation used here.  Applying this to \(\eta\)
and (A.154.32) proves
\(\eta=\kappa_2(Q)\) for a unique \(Q\in\{O,T_0\}\).

At the unique dyadic place, A.152 proves that the
local continuous Kummer map is injective and that the
proper fibres are exactly
\[
                  q_+^{-1}(O)=P_0+P_\infty,
             \qquad q_+^{-1}(T_0)=P_1+P_{1/2}.
\tag{A.154.33}
\]
Consequently an ordinary point of \(D^\circ(K_2)\)
cannot give the localization of \(\eta\).
The Selmer condition leaves precisely the cuspidal
possibilities in (A.154.4), proving the assertion.

\emph{The precise odd-place kernel.}
The dyadic conclusion cannot be propagated to every
odd place by the same Kummer comparison.  Indeed take
an odd place of good reduction for \(E_b\) and for the
displayed proper morphism.  Its formal subgroup
\(E_1(K_w)\) is uniquely two-divisible: in the formal
parameter multiplication by two has derivative two,
a local unit, and Hensel's inverse gives its inverse
on that subgroup.  Every element of \(E_1(K_w)\)
therefore has zero class in every finite two-power
Kummer quotient, and hence zero continuous Kummer
class by the preceding compatible-coboundary argument.
The map \(q_+\) is etale at the rational boundary point
\(P_0\) above \(O\), by (A.154.33).  The inverse
function form of Hensel's lemma gives an analytic
isomorphism between sufficiently small neighborhoods
of these two points.  Nonzero formal elliptic points
in that neighborhood lift to actual ordinary points
of \(D^\circ(K_w)\), while their two-primary Kummer
classes are zero.  These are the explicit maps and
points delimiting the all-place scope of this argument.
\end{proof}
